\chapter{Sintaxis \emph{chain} (legacy)}
\label{cap:dsl-chain}

\begin{intuicion}
La sintaxis \emph{chain} es la forma original de la DSL: una
cadena lineal de pasos \(p_0 \to p_1 \to \dots \to p_k\) con
predicados opcionales y cláusulas trailing. Es la forma más usada
y la única que existía antes de L1 stage 3c. Este capítulo
documenta cada parte de la sintaxis con ejemplos y muestra cómo el
parser la convierte en \texttt{Hypothesis}.
\end{intuicion}

\section{Anatomía de una hipótesis chain}

\begin{lstlisting}
User u -[Auth]-> Host h -[Connect]-> IP {k=2} HAVING count(distinct last) >= 5 WITHIN 30m
\end{lstlisting}

\noindent
Componentes:

\begin{itemize}
    \item \texttt{User u}: \emph{type declaration} con
    \emph{binding}. Declara que el primer vértice es de tipo
    \texttt{User} y se referenciará como \texttt{u}.
    \item \texttt{-[Auth]->}: arista con \texttt{RelationType}.
    \item \texttt{Host h}: type declaration del próximo vértice.
    \item \texttt{-[Connect]->}: segunda arista.
    \item \texttt{IP}: type declaration sin binding (no se
    referencia).
    \item \texttt{\{k=2\}}: \(k\)-simplicidad. Cada vértice puede
    repetirse hasta 2 veces.
    \item \texttt{HAVING count(distinct last) >= 5}: agregación
    post-hoc. Acepta sólo grupos de \(\geq 5\) lasts distintos.
    \item \texttt{WITHIN 30m}: ventana temporal de 30 minutos
    entre primera y última arista.
\end{itemize}

\section{Type bindings}

\begin{lstlisting}
User u -[Auth]-> Host h
\end{lstlisting}

\noindent
\texttt{u} y \texttt{h} son \emph{bindings}: identificadores que
referencian el vértice \emph{en cláusulas trailing} o \emph{en
otras chains} (shape graph). En una chain pura sin \texttt{HAVING},
los bindings son opcionales:

\begin{lstlisting}
User -[Auth]-> Host -[Connect]-> IP
\end{lstlisting}

\noindent
es equivalente a la versión con bindings, pero el motor pierde la
capacidad de referenciar los vértices intermedios.

\subsection{Casos donde el binding importa}

\begin{itemize}
    \item \texttt{HAVING count(distinct h) >= 3}: requiere binding
    \texttt{h} en el segundo vértice.
    \item Shape graph: \texttt{c -[Auth]-> Host h1; c -[Auth]-> Host h2}
    necesita binding \texttt{c} para que el segundo \texttt{c} lo
    referencie.
\end{itemize}

\section{Predicados sobre la arista}

\begin{lstlisting}
User -[Auth{result=success}]-> Host
\end{lstlisting}

\noindent
\texttt{\{result=success\}} es un predicado sobre la metadata de
la arista. El matcher consulta el \texttt{MetadataStore}
(capítulo~\ref{cap:entity-relation}) para verificar que el campo
\texttt{result} valga \texttt{success}.

Predicados soportados:

\begin{itemize}
    \item \texttt{key=value}: igualdad.
    \item \texttt{key>value}, \texttt{key<value}: comparación
    numérica.
    \item \texttt{key in [v1, v2, v3]}: pertenencia a conjunto.
    \item \texttt{key matches /regex/}: regex (caro, evitar si
    posible).
\end{itemize}

\section{Predicados sobre el vértice destino}

\begin{lstlisting}
User -[Auth]-> Host {os=Windows} -[Connect]-> IP
\end{lstlisting}

\noindent
Las llaves \emph{después} del type declaration aplican el predicado
al vértice destino. El matcher las verifica en el momento de
extender el path.

\section{Cláusulas trailing}

\subsection{\texttt{WITHIN}}

\begin{lstlisting}
User -[Auth]-> Host -[Connect]-> IP WITHIN 30m
\end{lstlisting}

\noindent
\(\Delta t = 30\) minutos entre el timestamp de la \emph{primera}
arista y el de la \emph{última}. Sintaxis: número + unidad
(\texttt{s}, \texttt{m}, \texttt{h}, \texttt{d}).

Sin \texttt{WITHIN}: \(\Delta t = +\infty\) (sin restricción).

\subsection{\texttt{HAVING}}

\begin{lstlisting}
User -[Auth]-> Host h -[Connect]-> IP HAVING count(distinct h) >= 5
\end{lstlisting}

\noindent
Agregación sobre los matches del path. \texttt{count(distinct h)}
cuenta hosts distintos por start; \texttt{>= 5} acepta sólo starts
con $\geq$ 5 hosts distintos.

Funciones soportadas:
\begin{itemize}
    \item \texttt{count(distinct binding)}.
    \item \texttt{count(distinct first)} / \texttt{count(distinct last)}:
    sintaxis especial referenciando el primer / último vértice de la
    chain.
    \item \texttt{count(*)}: cardinalidad bruta del grupo.
\end{itemize}

\subsection{\texttt{\{k=N\}}}

\begin{lstlisting}
User -[Auth]-> Host -[Connect]-> IP {k=2}
\end{lstlisting}

\noindent
\(k\)-simplicidad (capítulo~\ref{cap:grafos-temporales}): cada
vértice puede aparecer hasta \(k\) veces en el path. Default: \(k=1\)
(path simple). \(k=2\) permite revisitar un vértice una vez (útil
para hipótesis donde un usuario \emph{vuelve} a un host).

\section{Ejemplos completos}

\begin{ejemplo}[Lateral movement chain]
\begin{lstlisting}
User u -[Auth]-> Host h1 -[Auth]-> User u2 -[Auth]-> Host h2
WITHIN 1h
HAVING count(distinct h2) >= 3
\end{lstlisting}

\noindent
``Un usuario que se autentica en un host, desde donde otro usuario
se autentica en un tercer host distinto, repetido para al menos 3
hosts distintos, todo en menos de una hora.'' Patrón típico de
movimiento lateral con credential reuse.
\end{ejemplo}

\begin{ejemplo}[Beaconing inicial]
\begin{lstlisting}
Process p {image=powershell.exe} -[Connect]-> IP -[Beacon{interval<60s}]-> Domain
WITHIN 5m
\end{lstlisting}

\noindent
``Un proceso PowerShell que se conecta a una IP, y desde esa IP hay
beaconing con intervalo $<$ 60 segundos a un Domain, todo en menos
de 5 minutos.'' Detección C2 de bajo volumen.
\end{ejemplo}

\section{Cómo el parser construye la \texttt{Hypothesis}}

Para la hipótesis del primer ejemplo, el parser produce:

\begin{lstlisting}[language=Rust]
Hypothesis {
    name: "...",
    steps: vec![
        HypothesisStep {
            source_type: User,  source_binding: Some("u"),
            rel_type: Auth,     edge_predicates: vec![],
            dest_type: Host,    dest_binding: Some("h1"),
            dest_predicates: vec![],
        },
        HypothesisStep {
            source_type: Host,  source_binding: Some("h1"),
            rel_type: Auth,     edge_predicates: vec![],
            dest_type: User,    dest_binding: Some("u2"),
            dest_predicates: vec![],
        },
        HypothesisStep {
            source_type: User,  source_binding: Some("u2"),
            rel_type: Auth,     edge_predicates: vec![],
            dest_type: Host,    dest_binding: Some("h2"),
            dest_predicates: vec![],
        },
    ],
    k_simplicity: 1,
    aggregations: vec![
        AggClause::CountDistinct {
            binding: "h2",
            op: Op::Gte,
            value: 3,
        }
    ],
    window_seconds: Some(3600),
    shape: HypothesisShape::Chain,
}
\end{lstlisting}

\noindent
Notar que \texttt{shape == Chain} ---no hay \texttt{;} en la
hipótesis. Si lo hubiera, sería \texttt{Graph}
(capítulo~\ref{cap:dsl-graph}).

\section{Resumen}

\begin{itemize}
    \item Sintaxis chain: cadena lineal \texttt{Type binding -[rel]-> Type binding}
    + cláusulas trailing.
    \item Predicados: en aristas (\texttt{[rel\{predicates\}]}) y
    en vértices destino (\texttt{Type \{predicates\}}).
    \item Trailing: \texttt{WITHIN N(m|h|d)}, \texttt{HAVING agg op
    val}, \texttt{\{k=N\}}.
    \item Defaults: sin WITHIN $\Rightarrow$ ventana infinita;
    sin HAVING $\Rightarrow$ enumeración bruta; \texttt{\{k=1\}}
    $\Rightarrow$ path simple.
    \item \texttt{shape == Chain} para hipótesis sin \texttt{;}.
\end{itemize}

\begin{ejercicios}
\begin{enumerate}
    \item \starred{} Escribí en la DSL: ``un usuario que en menos
    de 10 minutos se autenticó en al menos 3 hosts distintos''.
    Identificá las cláusulas trailing necesarias.

    \item ¿Cuál es la diferencia entre
    \texttt{User -[Auth\{result=success\}]-> Host} y
    \texttt{User -[Auth]-> Host \{result=success\}}?

    \item \starred{} Construí la \texttt{Hypothesis} (struct Rust)
    para la hipótesis: \texttt{Process p -[Spawn]-> Process c -[Write]-> File f WITHIN 1m}.

    \item ¿Por qué los bindings son opcionales en chain pura pero
    obligatorios en HAVING? Pensá: ¿qué referenciaría el HAVING
    sin binding?

    \item Imaginá una hipótesis donde necesitás \(k=3\). Escribí
    un caso real (en threat hunting) donde la repetición triple es
    sensata.
\end{enumerate}
\end{ejercicios}
